\documentclass[11pt,a4paper]{article}

\usepackage[margin=1in]{geometry}
\usepackage{enumitem}
\usepackage{hyperref}
\usepackage{array}
\usepackage{booktabs}

\title{MiniChess Project Report}
\author{Student ID: 111061143}
\date{\today}

\begin{document}
\maketitle

\section{Overview}

This project implements a MiniChess state evaluator, legal move generation, and
a stronger benchmark policy. The main submitted files are
\texttt{111061143\_state.cpp}, \texttt{111061143\_state.hpp},
\texttt{111061143\_submission.cpp}, and \texttt{111061143\_submission.hpp}.
Additional integration files were also modified so the engine can register and
use the stronger search policy through UBGI.

\section{Modified Files}

\begin{center}
\begin{tabular}{>{\ttfamily}p{0.42\linewidth}p{0.48\linewidth}}
\toprule
File & Purpose \\
\midrule
src/games/minichess/111061143\_state.cpp &
Implements board evaluation and completes MiniChess legal move generation. \\
src/games/minichess/111061143\_state.hpp &
State interface used by the submitted implementation. \\
src/policy/111061143\_submission.cpp &
Implements MiniMax, Alpha-Beta, Principal Variation Search, move ordering,
quiescence search, and transposition table support. \\
src/policy/111061143\_submission.hpp &
Declares the search policies and search parameters. \\
src/policy/111061143\_registry.hpp &
Registers the available policies and changes the default policy to PVS. \\
src/ubgi/111061143\_ubgi.cpp &
Connects the submitted registry to the UBGI engine and adds a small White
opening book. \\
src/ubgi/111061143\_ubgi.hpp &
Includes the submitted state header. \\
Makefile &
Removes \texttt{-march=native} to improve compilation portability. \\
cli/cli.py &
Adds \texttt{alphabeta} and \texttt{pvs} to the command-line algorithm choices
and changes the default White/Black algorithm from \texttt{minimax} to
\texttt{pvs}. \\
\bottomrule
\end{tabular}
\end{center}

\section{State Implementation}

The state implementation completes the missing MiniChess logic in the original
starter code. In this project, a winning state means that a legal move can
capture the opponent king. Once this happens, the game is already decided, so
the evaluator returns the maximum score immediately instead of letting the
opponent move the king away on a later turn.
For non-terminal positions, the evaluator supports two modes:

\begin{itemize}[leftmargin=*]
    \item A simple material evaluator using piece values.
    \item A stronger evaluator using material, piece-square tables, king
    tropism, and mobility.
\end{itemize}

The stronger evaluator first locates both kings, then evaluates every piece from
the current player's perspective. Material counts the value of pieces on the
board. Piece-square tables give extra score for pieces on useful squares.
King tropism gives a smaller positional bonus when pieces are closer to the
opponent king; an immediate king capture is handled separately as a winning
state. Mobility is calculated by generating legal moves for both sides and
adding a bonus when the current player has more available moves.

The legal move generator was also completed for knights. The knight move table
contains the eight standard L-shaped moves, such as two squares in one
direction and one square sideways. For each candidate square, the code checks
whether the square is inside the board and whether it is already occupied by a
friendly piece. Legal empty-square moves and captures are added to the move
list. If one of those captures takes the opponent king, the state is marked as
a win and move generation returns immediately because the game is over.

\section{Search Policy}

The submitted policy keeps the original MiniMax interface but extends it with
stronger search algorithms. The engine still plays only one move in the current
turn, but it searches several future plies first. For example, depth 1 checks
the engine's next move, depth 2 also checks the opponent reply, and depth 3
checks the engine's next response after that. This lookahead helps avoid moves
that look good immediately but lose material or the king on the following turn.

\subsection{MiniMax}

The basic MiniMax implementation recursively evaluates child states to the
requested depth. Scores are converted back into the current player's
perspective after each child search. Terminal wins are scored as
\texttt{P\_MAX - ply}, which prefers faster wins. Draws and repeated positions
use a draw contempt parameter, making the engine slightly prefer continuing play
over accepting neutral repetition.

\subsection{Alpha-Beta Search}

Alpha-Beta pruning was added to reduce the number of explored nodes while
preserving the same final minimax value. Alpha is the best score the current
side has already found. Beta is the opponent's limit: if a branch becomes good
enough for one player that the other player would avoid allowing it, the engine
does not need to continue searching that branch. This pruning lets the engine
search deeper than plain MiniMax in the same amount of time.

\subsection{Principal Variation Search}

Principal Variation Search (PVS) is used as the default algorithm. PVS is an
optimized form of Alpha-Beta. It searches the first, best-looking move normally.
For later moves, it first performs a cheaper narrow-window search that only asks
whether the move can beat the current best score. If the answer is no, the move
is skipped quickly. If the answer is yes, the move is searched again with a full
window to get a more accurate score. This works well when move ordering puts a
strong move first.

\subsection{Move Ordering}

Moves are sorted before search. The ordering gives priority to:

\begin{itemize}[leftmargin=*]
    \item The best move from the transposition table.
    \item Captures, using a captured-piece-value minus moving-piece-value
    heuristic.
    \item Pawn promotions.
    \item Moves toward the center.
\end{itemize}

Good move ordering improves both Alpha-Beta and PVS because stronger moves are
searched earlier. Once a strong score is found, weaker branches can often be
cut off sooner, which saves time and allows deeper search.

\subsection{Transposition Table}

A thread-local transposition table stores previously searched positions by
state hash. This is useful because different move orders can sometimes reach
the same board position. Each entry records how deeply the position was
searched, the score, the type of score, and the best move found. An exact score
means the stored value is the searched value. A lower bound means the real score
is at least that good, and an upper bound means the real score is at most that
good. Cached results are only reused when the old search depth is at least as
deep as the current request, because a shallow result may miss future tactics.
The stored best move is also tried first next time, improving move ordering.

\subsection{Quiescence Search}

At depth zero, the search does not stop immediately on all positions. If the
position still has important tactical moves, the engine continues through noisy
moves such as captures and promotions. This reduces the horizon effect. For
example, if the engine stops while its queen can be captured immediately, it may
evaluate the position too optimistically. Quiescence search looks at the
immediate capture first, then evaluates a quieter and more stable position.

\section{UBGI Integration}

The registry adds three searchable policies: \texttt{minimax},
\texttt{alphabeta}, and \texttt{pvs}. The default algorithm is changed to
\texttt{pvs}, so benchmark runs use the strongest submitted search by default.

The UBGI file also includes a small exact-position book for White. It starts
from the opening position and may continue for many moves only if the board
keeps matching the stored line. The book is keyed by encoded board state and
side to move, so it is only used when the current position exactly matches a
stored position. If the opponent makes a different move, the position no longer
matches the book and the engine automatically falls back to normal PVS search.

\section{Parameters}

The submitted policy supports the following options:

\begin{center}
\begin{tabular}{>{\ttfamily}p{0.28\linewidth}p{0.16\linewidth}p{0.45\linewidth}}
\toprule
Parameter & Default & Description \\
\midrule
UseKPEval & true & Enables material, piece-square table, and king-tropism evaluation. \\
UseEvalMobility & true & Adds the legal-move mobility bonus. \\
ReportPartial & true & Reports root move updates during search. \\
DrawContempt & 20 & Penalizes draw and neutral repetition scores slightly. \\
\bottomrule
\end{tabular}
\end{center}

\section{Conclusion}

The final engine improves the starter MiniChess implementation in two main
ways. First, the game state is complete enough to evaluate positions and
generate knight moves correctly. Second, the submitted benchmark policy replaces
plain MiniMax with PVS, Alpha-Beta pruning, move ordering, transposition table
caching, and quiescence search. These changes increase search strength while
keeping the original project interfaces compatible with the provided benchmark
and UBGI framework.

\end{document}
